\section{Literature Review}\label{sec:literature}

The proposed problem connects two scheduling research streams. The first studies
completion-time costs defined by due dates or due windows and the exact solution of
the resulting machine-scheduling problems. The second integrates internal
scheduling with order rejection, acceptance, or outsourcing. We review each stream
and then identify the modelling and algorithmic gaps created by combining them.

\subsection{Scheduling with due-window costs}

Earliness--tardiness scheduling recognizes that completing every job as early as
possible is not generally optimal. The review of \citet{baker1990review} establishes
the classical role of inventory costs before a target date and service penalties
after it. Due-window models replace the target date by an interval of acceptable
completion times. \citet{wan2002tabu} study a single machine with distinct due
windows and weighted earliness--tardiness penalties. \citet{rosa2017algorithms}
consider distinct time windows, general early and late penalties, and
sequence-dependent setups; their complete problem is treated heuristically, while
exact results require a restricted setup structure. The annotated review of
\citet{kramer2019unified} places these variants within the broader family of
earliness--tardiness problems and documents the effects of machine environments,
release dates, setup assumptions, and permitted idle time.

Parallel-machine research has developed increasingly strong bounds and exact
decompositions for related completion-time objectives.
\citet{kedad2008lowerbounds} derive assignment-based and time-indexed lower bounds
for parallel-machine earliness--tardiness scheduling with distinct due dates.
\citet{sen2015preemptive} construct a preemptive relaxation for weighted tardiness
and earliness--tardiness problems on unrelated parallel machines and use it within
a Benders framework. \citet{bulhoes2020exact} propose a generic
branch-cut-and-price method in which each column is a complete schedule on one
machine. Their framework combines column generation with branching, cuts,
stabilization, and reduced-cost fixing for a broad class of parallel-machine
objectives.

These studies show both why due-window timing must be optimized and why a
single-machine schedule is a natural column. They do not decide whether a job
should leave the internal production system, nor do they account for a cost that
couples the jobs assigned to an external provider. The latter decision changes the
set of jobs entering every internal machine-pricing problem.

\subsection{Scheduling with outsourcing}

Order acceptance and scheduling jointly determine which orders enter a production
system and how accepted orders are processed. \citet{slotnick2011review} organize
this literature by machine environment, acceptance decision, and scheduling
criterion. Later models incorporate richer production structures. For example,
\citet{li2021acceptance} combine order acceptance with earliness--tardiness
penalties and overtime, while \citet{chen2023lbbd} apply logic-based Benders
decomposition to acceptance and scheduling across heterogeneous factories.
Rejection or nonacceptance is normally charged through an additive job-level loss,
so the cost of excluding one job does not depend on the other excluded jobs.

Scheduling with an explicit subcontracting option is more closely related to the
present decision. \citet{lee2008outsourcing} study single-machine scheduling with
outsourcing under due-date-related criteria. \citet{safarzadeh2019jobshop} allow
jobs or operations to be outsourced in a job-shop environment, and
\citet{kim2022flowshop} account for outsourcing lead times in a two-machine flow
shop. \citet{zhong2022coordinated} jointly optimize outsourcing, in-house parallel
machine production, and outbound delivery. Although these models differ in their
machine environments and objectives, their outsourcing payments are principally
specified job by job.

Aggregate pricing changes this separable structure. \citet{lu2021discount} study
in-house production and outsourcing under discount schemes defined on total
outsourcing cost. Their analysis establishes a scheduling precedent for pricing
the outsourced portfolio as a whole rather than summing independent adjusted job
charges. However, that setting does not include job-dependent due-window penalties,
sequence-dependent setups, or an exact column-generation treatment of internal
parallel-machine schedules.

\subsection{Synthesis and research gap}

The two streams leave a specific gap. Due-window scheduling captures the cost of
internal completion timing but keeps the job set fixed. Outsourcing scheduling
changes the internal job set but generally uses additive external costs or machine
environments without general due-window timing. We are not aware of an exact method
that combines identical parallel machines, job-dependent due windows,
sequence-dependent setup times and costs, and an aggregate quantity-discount
outsourcing tariff.

There is also an unresolved decomposition choice. Outsourcing can be represented by
one column for the complete outsourced set or by job-level decisions and tariff
variables in the master problem. These alternatives produce different master
relaxations, pricing families, and branching decisions. On the internal side,
existing exact parallel-machine methods motivate schedule-based column generation,
but time-indexed pricing expands with the numerical horizon. The formulations in
Section~\ref{sec:formulations} address the first issue, and the continuous-time
functional pricing method in Section~\ref{sec:solution} addresses the second.
